\chapter{2018年全国III卷（理）}

\section{单选题}

\begin{problem}
已知集合 \(A=\{x\mid x-1\ge0\}\)，\(B=\{0,1,2\}\)，则 \(A\cap B=\)（\quad）

\choices
    {\(\{0\}\)}
    {\(\{1\}\)}
    {\(\{1,2\}\)}
    {\(\{0,1,2\}\)}
\end{problem}

\begin{answer}
C
\end{answer}

\begin{solution}
由 \(x-1\ge0\) 得 \(x\ge1\)，所以 \(A\) 中属于 \(B=\{0,1,2\}\) 的元素为 \(1,2\)，即 \(A\cap B=\{1,2\}\)，故选 C.
\end{solution}

\begin{problem}
\((1+\i)(2-\i)=\)（\quad）

\choices
    {\(-3-\i\)}
    {\(-3+\i\)}
    {\(3-\i\)}
    {\(3+\i\)}
\end{problem}

\begin{answer}
D
\end{answer}

\begin{solution}
\[
(1+\i)(2-\i)=2-\i+2\i-\i^2=3+\i
\]
因此选 D.
\end{solution}

\begin{problem}
中国古建筑借助榫卯将木构件连接起来. 构件的凸出部分叫榫头，凹进部分叫卯眼，图中木构件右边的小长方体是榫头. 若如图摆放的木构件与某一带卯眼的木构件咬合成长方体，则咬合时带卯眼的木构件的俯视图可以是（\quad）

\bitmapfigure[max width=0.66\linewidth]{img/2018/national_paper_3_science/q3_fig1.png}

\choices
    {\choicebitmap[width=0.15\paperwidth]{img/2018/national_paper_3_science/q3_fig2.png}}
    {\choicebitmap[width=0.15\paperwidth]{img/2018/national_paper_3_science/q3_fig3.png}}
    {\choicebitmap[width=0.15\paperwidth]{img/2018/national_paper_3_science/q3_fig4.png}}
    {\choicebitmap[width=0.15\paperwidth]{img/2018/national_paper_3_science/q3_fig5.png}}
\end{problem}

\begin{answer}
A
\end{answer}

\begin{solution}
从俯视方向看，榫头的投影是大木构件右侧中部的矩形，因此带卯眼的木构件必须在左侧中部有一个与榫头等大、等高的矩形凹槽；凹槽的内边界被上方木板遮挡，俯视图中应画成虚线，同时木构件的外轮廓仍应是完整矩形. 第一项恰好满足这三个条件. 第二项把凹槽画成外露的实线内框，第三项的凹槽未与左侧边缘连通，第四项的凹槽直接破坏了外轮廓，均不能与榫头完全咬合后组成完整长方体，故选 A.
\end{solution}

\begin{problem}
若 \(\sin\alpha=\dfrac13\)，则 \(\cos2\alpha=\)（\quad）

\choices
    {\(\dfrac89\)}
    {\(\dfrac79\)}
    {\(-\dfrac79\)}
    {\(-\dfrac89\)}
\end{problem}

\begin{answer}
B
\end{answer}

\begin{solution}
由二倍角公式，
\[
\cos 2\alpha=1-2\sin^2\alpha=1-2\left(\dfrac13\right)^2=\dfrac79
\]
因此选 B.
\end{solution}

\begin{problem}
\(\left(x^2+\dfrac{2}{x}\right)^5\) 的展开式中 \(x^4\) 的系数为（\quad）

\choices
    {\(10\)}
    {\(20\)}
    {\(40\)}
    {\(80\)}
\end{problem}

\begin{answer}
C
\end{answer}

\begin{solution}
在展开式中，设取了 \(k\) 个因式 \(\dfrac2x\)，则相应项为
\[
\mathrm C_5^k(x^2)^{5-k}\left(\dfrac2x\right)^k
=\mathrm C_5^k2^k x^{10-3k}
\]
要得到 \(x^4\)，需 \(10-3k=4\)，从而 \(k=2\). 所以 \(x^4\) 的系数为
\[
\mathrm C_5^2 2^2=10\times4=40
\]
故选 C.
\end{solution}

\begin{problem}
直线 \(x+y+2=0\) 分别与 \(x\) 轴，\(y\) 轴交于 \(A\)，\(B\) 两点，点 \(P\) 在圆 \((x-2)^2+y^2=2\) 上，则 \(\triangle ABP\) 面积的取值范围是（\quad）

\choices
    {\([2,6]\)}
    {\([4,8]\)}
    {\(\left[\sqrt2,3\sqrt2\right]\)}
    {\(\left[2\sqrt2,3\sqrt2\right]\)}
\end{problem}

\begin{answer}
A
\end{answer}

\begin{solution}
直线 \(x+y+2=0\) 与坐标轴的交点为 \(A(-2,0)\)，\(B(0,-2)\)，故
\[
|AB|=\sqrt{(2)^2+(-2)^2}=2\sqrt2
\]
点 \(P(x,y)\) 到直线 \(AB\) 的距离为
\[
\dfrac{|x+y+2|}{\sqrt2}
\]
因此
\[
S_{\triangle ABP}=\dfrac12|AB|\cdot\dfrac{|x+y+2|}{\sqrt2}=|x+y+2|
\]
圆心为 \(O_1(2,0)\)，半径为 \(\sqrt2\). 线性函数 \(x+y+2\) 在该圆上的中心值为 \(4\)，波动范围为
\[
\sqrt2\sqrt{1^2+1^2}=2
\]
所以 \(x+y+2\in[2,6]\)，从而三角形面积的取值范围为 \([2,6]\)，故选 A.
\end{solution}

\begin{problem}
函数 \(y=-x^4+x^2+2\) 的图象大致为（\quad）

\choices
    {\choicebitmap[width=0.15\paperwidth]{img/2018/national_paper_3_science/q7_fig1.png}}
    {\choicebitmap[width=0.15\paperwidth]{img/2018/national_paper_3_science/q7_fig2.png}}
    {\choicebitmap[width=0.15\paperwidth]{img/2018/national_paper_3_science/q7_fig3.png}}
    {\choicebitmap[width=0.15\paperwidth]{img/2018/national_paper_3_science/q7_fig4.png}}
\end{problem}

\begin{answer}
D
\end{answer}

\begin{solution}
函数 \(f(x)=-x^4+x^2+2\) 是偶函数，且 \(f(0)=2\)，\(f(\pm1)=2\). 其导数为
\[
f'(x)=-4x^3+2x=2x(1-2x^2)
\]
所以极值点为 \(x=0\) 和 \(x=\pm\dfrac1{\sqrt2}\)，其中
\(f(0)=2\)，\(f\left(\pm\dfrac1{\sqrt2}\right)=\dfrac94\).
又由 \(-x^4+x^2+2=0\) 得 \(x=\pm\sqrt2\)，且当 \(|x|\to\infty\) 时 \(f(x)\to-\infty\). 因此图象关于 \(y\) 轴对称，在原点附近有一个较低的局部极小值，在 \(x=\pm\dfrac1{\sqrt2}\) 处有两个局部极大值，并在 \(x=\pm\sqrt2\) 处与 \(x\) 轴相交，符合选项 D，故选 D.
\end{solution}

\begin{problem}
某群体中的每位成员使用移动支付的概率都为 \(p\)，各成员的支付方式相互独立. 设 \(X\) 为该群体的 \(10\) 位成员中使用移动支付的人数，\(D(X)=2.4\)，\(P(X=4)<P(X=6)\)，则 \(p=\)（\quad）

\choices
    {\(0.7\)}
    {\(0.6\)}
    {\(0.4\)}
    {\(0.3\)}
\end{problem}

\begin{answer}
B
\end{answer}

\begin{solution}
由 \(X\sim B(10,p)\)，得
\[
D(X)=10p(1-p)=2.4
\]
即 \(p(1-p)=0.24\)，所以 \(p=0.4\) 或 \(p=0.6\). 又
\[
\dfrac{P(X=4)}{P(X=6)}
=\dfrac{\mathrm C_{10}^4p^4(1-p)^6}{\mathrm C_{10}^6p^6(1-p)^4}
=\left(\dfrac{1-p}{p}\right)^2
\]
题设 \(P(X=4)<P(X=6)\)，故 \((1-p)/p<1\)，从而 \(p>\dfrac12\). 因此 \(p=0.6\)，故选 B.
\end{solution}

\begin{problem}
\(\triangle ABC\) 的内角 \(A\)，\(B\)，\(C\) 的对边分别为 \(a\)，\(b\)，\(c\). 若 \(\triangle ABC\) 的面积为 \(\dfrac{a^2+b^2-c^2}{4}\)，则 \(C=\)（\quad）

\choices
    {\(\dfrac{\pi}{6}\)}
    {\(\dfrac{\pi}{3}\)}
    {\(\dfrac{\pi}{4}\)}
    {\(\dfrac{\pi}{2}\)}
\end{problem}

\begin{answer}
C
\end{answer}

\begin{solution}
三角形面积为
\[
S_{\triangle ABC}=\dfrac12ab\sin C
\]
由余弦定理，\(a^2+b^2-c^2=2ab\cos C\)，所以题设等式化为
\[
\dfrac12ab\sin C=\dfrac12ab\cos C
\]
因为 \(a\)，\(b>0\)，且面积为正，故 \(\sin C=\cos C>0\)，从而 \(C=\dfrac\pi4\)，故选 C.
\end{solution}

\begin{problem}
设 \(A\)，\(B\)，\(C\)，\(D\) 是同一个半径为 \(4\) 的球的球面上四点，\(\triangle ABC\) 为等边三角形且面积为 \(9\sqrt3\)，则三棱锥 \(D-ABC\) 体积的最大值为（\quad）

\choices
    {\(12\sqrt3\)}
    {\(18\sqrt3\)}
    {\(24\sqrt3\)}
    {\(54\sqrt3\)}
\end{problem}

\begin{answer}
B
\end{answer}

\begin{solution}
设等边三角形 \(ABC\) 的边长为 \(s\). 由面积条件，
\(\dfrac{\sqrt3}{4}s^2=9\sqrt3\)，\(s=6\).
其外接圆半径为 \(\dfrac{s}{\sqrt3}=2\sqrt3\). 设球心为 \(O\)，\(O\) 到平面 \(ABC\) 的距离为 \(h\)，则
\[
h=\sqrt{4^2-(2\sqrt3)^2}=2
\]
三棱锥 \(D-ABC\) 的体积为
\[
V=\dfrac13S_{\triangle ABC}\cdot d(D,\text{平面 }ABC)
\]
点 \(D\) 在半径为 \(4\) 的球面上，故其到平面 \(ABC\) 的距离最大为 \(4+2=6\). 于是
\[
V_{\max}=\dfrac13\cdot9\sqrt3\cdot6=18\sqrt3
\]
故选 B.
\end{solution}

\begin{problem}
设 \(F_1\)，\(F_2\) 是双曲线 \(C:\dfrac{x^2}{a^2}-\dfrac{y^2}{b^2}=1\)（\(a>0\)，\(b>0\)）的左，右焦点，\(O\) 是坐标原点. 过 \(F_2\) 作 \(C\) 的一条渐近线的垂线，垂足为 \(P\)，若 \(\left|PF_1\right|=\sqrt6|OP|\)，则 \(C\) 的离心率为（\quad）

\choices
    {\(\sqrt5\)}
    {\(2\)}
    {\(\sqrt3\)}
    {\(\sqrt2\)}
\end{problem}

\begin{answer}
C
\end{answer}

\begin{solution}
设双曲线的焦距参数为 \(c\)，则 \(c^2=a^2+b^2\)，且 \(F_2=(c,0)\). 取渐近线 \(y=\dfrac ba x\)，其方向向量可取 \((a,b)\)，长度为 \(c\). 从 \(F_2\) 到该渐近线的垂足为
\[
P=\dfrac{(c,0)\cdot(a,b)}{c^2}(a,b)=\left(\dfrac{a^2}{c},\dfrac{ab}{c}\right)
\]
因此 \(|OP|=a\). 又 \(F_1=(-c,0)\)，所以
\[
|PF_1|^2=\left(\dfrac{a^2}{c}+c\right)^2+\left(\dfrac{ab}{c}\right)^2
\]
利用 \(c^2=a^2+b^2\)，整理得
\[
|PF_1|^2=\dfrac{4a^4+5a^2b^2+b^4}{c^2}
\]
由题设 \(|PF_1|=\sqrt6|OP|=\sqrt6a\)，故
\[
4a^4+5a^2b^2+b^4=6a^2c^2=6a^2(a^2+b^2)
\]
令 \(t=\dfrac{b^2}{a^2}>0\)，约去 \(a^4\) 得
\(4+5t+t^2=6+6t\)，\(t^2-t-2=0\).
所以 \(t=2\)，即 \(b^2=2a^2\). 双曲线离心率为
\[
e=\dfrac ca=\sqrt{1+\dfrac{b^2}{a^2}}=\sqrt3
\]
故选 C.
\end{solution}

\begin{problem}
设 \(a=\log_{0.2}0.3\)，\(b=\log_20.3\)，则（\quad）

\choices
    {\(a+b<ab<0\)}
    {\(ab<a+b<0\)}
    {\(a+b<0<ab\)}
    {\(ab<0<a+b\)}
\end{problem}

\begin{answer}
B
\end{answer}

\begin{solution}
因为 \(0.2<0.3<1\)，所以 \(0<a=\log_{0.2}0.3<1\)；又因为 \(2>1\) 且 \(0<0.3<1\)，所以 \(b=\log_2 0.3<0\). 设
\(t=-\ln0.3>0\)，\(u=\ln5\)，\(v=\ln2\).
则 \(a=\dfrac tu\)，\(b=-\dfrac tv\). 由于 \(u>v\)，有
\(a+b=t\left(\dfrac1u-\dfrac1v\right)<0\)，\(ab<0\).
此外，\(ab<a+b\) 等价于
\[
-\dfrac{t^2}{uv}<\dfrac{t(v-u)}{uv}
\]
即 \(u-v<t\). 而 \(u-v=\ln\dfrac52\)，\(t=\ln\dfrac{10}{3}\)，且 \(\dfrac52<\dfrac{10}{3}\)，故 \(u-v<t\). 因此 \(ab<a+b<0\)，故选 B.
\end{solution}

\section{填空题}

\begin{problem}
已知向量 \(\bs a=(1,2)\)，\(\bs b=(2,-2)\)，\(\bs c=(1,\lambda)\). 若 \(\bs c\parallel (2\bs a+\bs b)\)，则 \(\lambda=\fillinblank{}\).
\end{problem}

\begin{answer}
\(\dfrac12\)
\end{answer}

\begin{solution}
先算得
\[
2\bs a+\bs b=2(1,2)+(2,-2)=(4,2)
\]
由于 \(\bs c=(1,\lambda)\) 与 \((4,2)\) 平行，故对应分量成比例：
\[
\dfrac{1}{4}=\dfrac{\lambda}{2}
\]
从而 \(\lambda=\dfrac12\).
\end{solution}

\begin{problem}
曲线 \(y=(ax+1)\e^x\) 在点 \((0,1)\) 处的切线的斜率为 \(-2\)，则 \(a=\fillinblank{}\).
\end{problem}

\begin{answer}
\(-3\)
\end{answer}

\begin{solution}
对 \(y=(ax+1)\e^x\) 求导，得
\[
y'=a\e^x+(ax+1)\e^x=(ax+a+1)\e^x
\]
在点 \((0,1)\) 处切线斜率为 \(y'(0)=a+1\). 由题设 \(a+1=-2\)，所以 \(a=-3\).
\end{solution}

\begin{problem}
函数 \(f(x)=\cos\!\left(3x+\dfrac{\pi}{6}\right)\) 在 \([0,\pi]\) 的零点个数为\fillinblank{}.
\end{problem}

\begin{answer}
\(3\)
\end{answer}

\begin{solution}
令 \(\cos\left(3x+\dfrac\pi6\right)=0\)，则
\(3x+\dfrac\pi6=\dfrac\pi2+k\pi\)，\(k\in\Z\)
从而
\[
x=\dfrac\pi9+\dfrac{k\pi}{3}
\]
当 \(x\in[0,\pi]\) 时，\(k=0\)，\(1\)，\(2\)，对应三个零点
\(\dfrac\pi9\)，\(\dfrac{4\pi}{9}\)，\(\dfrac{7\pi}{9}\).
所以零点个数为 \(3\).
\end{solution}

\begin{problem}
已知点 \(M(-1,1)\) 和抛物线 \(C:y^2=4x\)，过 \(C\) 的焦点且斜率为 \(k\) 的直线与 \(C\) 交于 \(A\)，\(B\) 两点. 若 \(\angle AMB=90^\circ\)，则 \(k=\fillinblank{}\).
\end{problem}

\begin{answer}
\(2\)
\end{answer}

\begin{solution}
抛物线 \(y^2=4x\) 的焦点为 \(F(1,0)\). 设直线与抛物线交于参数点
\(A=(t_1^2,2t_1)\)，\(B=(t_2^2,2t_2)\).
由于直线过焦点且斜率为 \(k\)，有 \(2t=k(t^2-1)\)，故 \(t_1\)，\(t_2\) 是方程
\[
kt^2-2t-k=0
\]
的两个根. 因为该直线与抛物线有两个交点，\(k\ne0\)，所以由韦达定理得 \(t_1t_2=-1\)，并且 \(t_1+t_2=\dfrac2k\).

由 \(\angle AMB=90^\circ\)，有
\[
\overrightarrow{MA}\cdot\overrightarrow{MB}=0
\]
因 \(M=(-1,1)\)，所以
\(\overrightarrow{MA}=(t_1^2+1,2t_1-1)\)，\(\overrightarrow{MB}=(t_2^2+1,2t_2-1)\).
利用 \(t_1t_2=-1\)，令 \(s=t_1+t_2\)，则
\[
\overrightarrow{MA}\cdot\overrightarrow{MB}
=(t_1^2+1)(t_2^2+1)+(2t_1-1)(2t_2-1)
=s^2+4-3-2s=(s-1)^2
\]
故 \(s=1\)，即 \(\dfrac2k=1\)，所以 \(k=2\).
\end{solution}

\section{解答题}

\begin{problem}
等比数列 \(\{a_n\}\) 中，\(a_1=1\)，\(a_5=4a_3\).

\begin{enumerate}
\item 求 \(\{a_n\}\) 的通项公式；
\item 记 \(S_n\) 为 \(\{a_n\}\) 的前 \(n\) 项和. 若 \(S_m=63\)，求 \(m\).
\end{enumerate}
\end{problem}

\begin{solution}
\begin{enumerate}
\item 设等比数列的公比为 \(q\). 由 \(a_1=1\) 得 \(a_3=q^2\)，\(a_5=q^4\)，代入 \(a_5=4a_3\) 得
\[
q^4=4q^2
\]
由 \(q^2(q^2-4)=0\)，代数上得到 \(q=0\)，\(q=2\) 或 \(q=-2\). 按通常把等比数列公比取为非零数的约定，\(q=0\) 不作为等比数列的公比；即使将它作为递推关系的退化情形一并考察，也有
\(a_1=1\)，\(a_n=0\)，\((n\ge2)\).
当 \(q=2\) 或 \(q=-2\) 时，通项公式分别为
\(a_n=2^{n-1}\) 或 \(a_n=(-2)^{n-1}\).
\item 退化情形 \(q=0\) 给出 \(S_m=1\)，不满足 \(S_m=63\). 当 \(q=2\) 时，
\[
S_m=1+2+\cdots+2^{m-1}=2^m-1
\]
由 \(S_m=63\) 得 \(2^m=64\)，所以 \(m=6\). 当 \(q=-2\) 时，
\[
S_m=\dfrac{1-(-2)^m}{1-(-2)}=\dfrac{1-(-2)^m}{3}
\]
若 \(S_m=63\)，则 \((-2)^m=-188\)，这不可能. 因此 \(m=6\).
\end{enumerate}
\end{solution}

\begin{problem}
某工厂为提高生产效率，开展技术创新活动，提出了完成某项生产任务的两种新的生产方式. 为比较两种生产方式的效率，选取 \(40\) 名工人，将他们随机分成两组，每组 \(20\) 人. 第一组工人用第一种生产方式，第二组工人用第二种生产方式. 根据工人完成生产任务的工作时间（单位：\(\mathrm{min}\)）绘制了如图茎叶图：

\begin{center}
\begin{tabular}{r|c|l}
第一种生产方式 & 茎 & 第二种生产方式 \\
\hline
\(8\) & \(6\) & \(5\)，\(5\)，\(6\)，\(8\)，\(9\) \\
\(9\)，\(7\)，\(6\)，\(2\) & \(7\) & \(0\)，\(1\)，\(2\)，\(2\)，\(3\)，\(4\)，\(5\)，\(6\)，\(6\)，\(8\) \\
\(9\)，\(8\)，\(7\)，\(7\)，\(6\)，\(5\)，\(4\)，\(3\)，\(3\)，\(2\) & \(8\) & \(1\)，\(4\)，\(4\)，\(5\) \\
\(2\)，\(1\)，\(1\)，\(0\)，\(0\) & \(9\) & \(0\)
\end{tabular}
\end{center}

\begin{enumerate}
\item 根据茎叶图判断哪种生产方式的效率更高？并说明理由；
\item 求 \(40\) 名工人完成生产任务所需时间的中位数 \(m\)，并将完成生产任务所需时间超过 \(m\) 和不超过 \(m\) 的工人数填入下面的列联表：
\begin{center}
\begin{tabular}{c|c|c}
 & 超过 \(m\) & 不超过 \(m\) \\
\hline
第一种生产方式 &  &  \\
\hline
第二种生产方式 &  &  \\
\end{tabular}
\end{center}

\item 根据（2）中的列联表，能否有 \(99\%\) 的把握认为两种生产方式的效率有差异？
\end{enumerate}

附：\(K^2=\dfrac{n\left( ad-bc \right)^2}{\left( a+b \right)\left( c+d \right)\left( a+c \right)\left( b+d \right)}\)，

\begin{center}
\begin{tabular}{c|c|c|c}
\(P\left( K^2\ge k \right)\) & 0.050 & 0.010 & 0.001 \\
\hline
\(k\) & 3.841 & 6.635 & 10.828
\end{tabular}
\end{center}
\end{problem}

\begin{solution}
\begin{enumerate}
\item 第一种生产方式的完成时间主要集中在 \(80\) 分钟以上，而第二种生产方式的完成时间主要集中在 \(80\) 分钟以下. 例如，第一种生产方式的中位数为 \(\dfrac{85+86}{2}=85.5\)，第二种生产方式的中位数为 \(\dfrac{73+74}{2}=73.5\). 完成同一任务所需时间越短，效率越高，因此第二种生产方式的效率更高.
\item 将两组数据合并并按从小到大排列，排在中间的两个数据为 \(79\) 和 \(81\)，故
\[
m=\dfrac{79+81}{2}=80
\]
第一种生产方式的数据中不超过 \(m\) 分钟的有 \(5\) 人，超过 \(m\) 分钟的有 \(15\) 人；第二种生产方式的数据中不超过 \(m\) 分钟的有 \(15\) 人，超过 \(m\) 分钟的有 \(5\) 人. 因此列联表为
\begin{center}
\begin{tabular}{c|c|c}
 & 超过 \(m\) & 不超过 \(m\) \\
\hline
第一种生产方式 & \(15\) & \(5\) \\
\hline
第二种生产方式 & \(5\) & \(15\)
\end{tabular}
\end{center}
\item 令列联表中的 \(a\)，\(b\)，\(c\)，\(d\) 依次为 \(15\)，\(5\)，\(5\)，\(15\)，则
\[
K^2=\dfrac{40(15\times15-5\times5)^2}{(15+5)(5+15)(15+5)(5+15)}
=10
\]
由于 \(10>6.635\)，故有 \(99\%\) 的把握认为两种生产方式的效率有差异.
\end{enumerate}
\end{solution}

\begin{problem}
如图，矩形 \(ABCD\) 所在的平面与半圆弧 \(\widehat{CD}\) 所在平面垂直，\(M\) 是 \(\widehat{CD}\) 上异于 \(C\)，\(D\) 的点.

\begin{enumerate}
\item 证明：平面 \(AMD\perp\) 平面 \(BMC\)；
\item 当三棱锥 \(M-ABC\) 体积最大时，求面 \(MAB\) 与面 \(MCD\) 所成二面角的正弦值.
\end{enumerate}

\bitmapfigure[max width=0.42\linewidth]{img/2018/national_paper_3_science/q19_fig1.png}
\end{problem}

\begin{solution}
\begin{enumerate}
\item 设 \(AB=a\)，\(AD=b\)，其中 \(a\)，\(b>0\). 建立空间直角坐标系，使
\(A(0,0,0)\)，\(B(a,0,0)\)，\(C(a,b,0)\)，\(D(0,b,0)\).
半圆弧 \(\widehat{CD}\) 所在平面为 \(y=b\). 设 \(M=(u,b,v)\)，其中 \(v>0\). 由 \(CD\) 为该半圆的直径，得
\[
\left(u-\dfrac a2\right)^2+v^2=\dfrac{a^2}{4}
\]
即 \(u^2-au+v^2=0\). 平面 \(AMD\) 和平面 \(BMC\) 的法向量分别可取
\(\bs n_1=\overrightarrow{AD}\times\overrightarrow{DM}=(bv,0,-bu)\)，\(\bs n_2=\overrightarrow{BC}\times\overrightarrow{CM}=(bv,0,b(a-u))\).
于是
\[
\bs n_1\cdot\bs n_2=b^2\left(v^2-au+u^2\right)=0
\]
所以 \(\bs n_1\perp\bs n_2\)，从而平面 \(AMD\perp\) 平面 \(BMC\).
\item 三棱锥 \(M-ABC\) 的底面面积为 \(S_{\triangle ABC}=\dfrac12ab\)，其高为 \(v\)，所以体积
\[
V=\dfrac16abv
\]
由 \(\left(u-\dfrac a2\right)^2+v^2=\dfrac{a^2}{4}\) 得 \(v\le\dfrac a2\)，当 \(M\) 为半圆弧中点，即 \(M=\left(\dfrac a2,b,\dfrac a2\right)\) 时，体积取得最大值.

此时平面 \(MCD\) 为 \(y=b\)，其法向量可取 \(\bs n_2=(0,1,0)\). 平面 \(MAB\) 的两个方向向量可取
\(\overrightarrow{AB}=(a,0,0)\)，\(\overrightarrow{AM}=\left(\dfrac a2,b,\dfrac a2\right)\)
故其法向量可取
\[
\bs n_1=\overrightarrow{AB}\times\overrightarrow{AM}=\left(0,-\dfrac{a^2}{2},ab\right)
\]
设两平面的二面角为 \(\varphi\)，则
\[
\sin\varphi=\dfrac{|\bs n_1\times\bs n_2|}{|\bs n_1||\bs n_2|}
=\dfrac{ab}{a\sqrt{\dfrac{a^2}{4}+b^2}}
=\dfrac{b}{\sqrt{\dfrac{a^2}{4}+b^2}}
\]
该值随矩形的边长比改变. 例如 \(a=b\) 时为 \(\dfrac2{\sqrt5}\)，而 \(a=2b\) 时为 \(\dfrac1{\sqrt2}\). 题目只给出“矩形”而未给出 \(a\)、\(b\) 或其比值，因此第（2）问不能确定唯一数值；这里保留由原文件题面能够推出的参数表达式，不以公开版本中的其他附加条件替换原题.
\end{enumerate}
\end{solution}

\begin{problem}
已知斜率为 \(k\) 的直线 \(l\) 与椭圆 \(C:\dfrac{x^2}{4}+\dfrac{y^2}{3}=1\) 交于 \(A\)，\(B\) 两点，线段 \(AB\) 的中点为 \(M(1,m)\)（\(m>0\)）.

\begin{enumerate}
\item 证明：\(k<-\dfrac12\)；
\item 设 \(F\) 为 \(C\) 的右焦点，\(P\) 为 \(C\) 上一点，且 \(\overrightarrow{FP}+\overrightarrow{FA}+\overrightarrow{FB}=\overrightarrow{0}\). 证明：\(\left|\overrightarrow{FA}\right|\)，\(\left|\overrightarrow{FP}\right|\)，\(\left|\overrightarrow{FB}\right|\) 成等差数列，并求该数列的公差.
\end{enumerate}
\end{problem}

\begin{solution}
\begin{enumerate}
\item 设 \(A(x_1,y_1)\)，\(B(x_2,y_2)\). 过中点 \(M(1,m)\) 的直线方程为
\[
y-m=k(x-1)
\]
将其代入椭圆方程，得到关于 \(x\) 的二次方程
\[
3x^2+4\bigl(kx+m-k\bigr)^2=12
\]
由韦达定理，
\[
x_1+x_2=\dfrac{-8k(m-k)}{3+4k^2}
\]
因 \(M\) 是 \(AB\) 的中点，\(x_1+x_2=2\)，所以
\(-8k(m-k)=2(3+4k^2)\)，\(km=-\dfrac34\).
又因为 \(M\) 是椭圆内弦的中点，故
\[
\dfrac{1^2}{4}+\dfrac{m^2}{3}<1
\]
从而 \(0<m<\dfrac32\). 由 \(km=-\dfrac34\) 得
\[
k=-\dfrac{3}{4m}< -\dfrac12
\]
\item 由向量条件和中点关系，
\[
\overrightarrow{FA}+\overrightarrow{FB}=2\overrightarrow{FM}=(0,2m)
\]
因此
\[
\overrightarrow{FP}=-(0,2m)=(0,-2m)
\]
椭圆的焦半距为 \(c=\sqrt{4-3}=1\)，故右焦点为 \(F=(1,0)\)，从而 \(P=(1,-2m)\). 因为 \(P\) 在椭圆上，
\[
\dfrac{1^2}{4}+\dfrac{(-2m)^2}{3}=1
\]
所以 \(m=\dfrac34\)，从而 \(k=-1\)，且 \(P=(1,-\dfrac32\)).

令 \(A=(1+u,\dfrac34-u)\)，\(B=(1-u,\dfrac34+u)\). 将 \(A\) 代入椭圆方程，得
\[
\dfrac{(1+u)^2}{4}+\dfrac{(\dfrac34-u)^2}{3}=1
\]
即 \(u^2=\dfrac{27}{28}\)，所以 \(|u|=\dfrac{3\sqrt{21}}{14}\). 对于椭圆上任一点 \(X(x,y)\)，由椭圆方程得 \(y^2=3-\dfrac34x^2\)，所以
\[
XF^2=(x-1)^2+y^2=x^2-2x+1+3-\dfrac34x^2=\dfrac{(4-x)^2}{4}
\]
又因椭圆上 \(x\in[-2,2]\)，所以 \(XF=2-\dfrac x2\). 于是
\(|FA|=\dfrac32-\dfrac u2\)，\(|FP|=\dfrac32\)，\(|FB|=\dfrac32+\dfrac u2\)
其中按 \(A\)，\(B\) 的顺序取 \(u\) 或 \(-u\) 只会交换首末两项. 因此三者成等差数列，公差的绝对值为
\[
\dfrac{|u|}{2}=\dfrac{3\sqrt{21}}{28}
\]
由于 \(A\)，\(B\) 的标号可互换，按题中顺序取公差时
\[
d=\pm\dfrac{3\sqrt{21}}{28}
\]
\end{enumerate}
\end{solution}

\begin{problem}
已知函数 \(f(x)=(2+x+ax^2)\ln(1+x)-2x\).

\begin{enumerate}
\item 若 \(a=0\)，证明：当 \(-1<x<0\) 时，\(f(x)<0\)；当 \(x>0\) 时，\(f(x)>0\)；
\item 若 \(x=0\) 是 \(f(x)\) 的极大值点，求 \(a\).
\end{enumerate}
\end{problem}

\begin{solution}
\begin{enumerate}
\item 当 \(a=0\) 时，\(f(x)=(x+2)\ln(1+x)-2x\)，定义域为 \(x>-1\). 求导得
\[
f'(x)=\ln(1+x)-\dfrac{x}{1+x}
\]
令 \(t=1+x>0\)，则
\[
f'(x)=\ln t-1+\dfrac1t
\]
记 \(g(t)=\ln t-1+\dfrac1t\)，则
\[
g'(t)=\dfrac{t-1}{t^2}
\]
所以 \(g(t)\) 在 \(0<t<1\) 上递减，在 \(t>1\) 上递增，且 \(g(1)=0\)，从而 \(g(t)\ge0\). 故 \(f(x)\) 在 \((-1,+\infty)\) 上递增. 又 \(f(0)=0\)，所以当 \(-1<x<0\) 时 \(f(x)<0\)，当 \(x>0\) 时 \(f(x)>0\).
\item 在 \(x=0\) 附近使用
\[
\ln(1+x)=x-\dfrac{x^2}{2}+\dfrac{x^3}{3}-\dfrac{x^4}{4}+o(x^4)
\]
将其代入 \(f(x)=(2+x+ax^2)\ln(1+x)-2x\) 并整理，得
\[
f(x)=\left(a+\dfrac16\right)x^3-\left(\dfrac16+\dfrac a2\right)x^4+o(x^4)
\]
若 \(a+\dfrac16\ne0\)，则首项为奇次项，\(f(x)\) 在 \(0\) 的两侧符号相反，\(x=0\) 不可能是极大值点. 因此必须有 \(a=-\dfrac16\). 此时
\[
f(x)=-\dfrac1{12}x^4+o(x^4)<0
\]
对充分接近 \(0\) 且不等于 \(0\) 的 \(x\) 成立，所以 \(x=0\) 确为极大值点.
\end{enumerate}
\end{solution}

\section{选考题：解答题}

\begin{problem}
在平面直角坐标系 \(xOy\) 中，\(\odot O\) 的参数方程为
\examdisplaycases{}{
x=\cos\theta,\\
y=\sin\theta
}
\(\theta\) 为参数，
过点 \((0,-\sqrt2)\) 且倾斜角为 \(\alpha\) 的直线 \(l\) 与 \(\odot O\) 交于 \(A\)，\(B\) 两点.

\begin{enumerate}
\item 求 \(\alpha\) 的取值范围；
\item 求 \(AB\) 中点 \(P\) 的轨迹的参数方程.
\end{enumerate}
\end{problem}

\begin{solution}
\begin{enumerate}
\item 取直线 \(l\) 的单位方向向量为 \((\cos\alpha,\sin\alpha)\). 点 \((0,-\sqrt2)\) 到 \(l\) 的距离为
\[
\sqrt2|\cos\alpha|
\]
直线与单位圆有两个交点的充要条件是该距离小于半径 \(1\)，故
\[
\sqrt2|\cos\alpha|<1
\]
在倾斜角 \(0\le\alpha<\pi\) 的约定下，得到
\[
\alpha\in\left(\dfrac\pi4,\dfrac{3\pi}4\right)
\]
\item 设 \(P=(x,y)\). 直线 \(l\) 的方程可写为
\[
-\sin\alpha\,x+\cos\alpha\,(y+\sqrt2)=0
\]
弦 \(AB\) 的中点 \(P\) 是圆心 \(O\) 到直线 \(l\) 的垂足，因此由点到直线的垂足公式，
\(x=\sqrt2\sin\alpha\cos\alpha\)，\(y=-\sqrt2\cos^2\alpha\).
这已经给出了轨迹的参数方程，参数为 \(\alpha\in\left(\dfrac\pi4,\dfrac{3\pi}4\right)\). 若令 \(m=\cot\alpha\)，则 \(-1<m<1\)，并且
\(x=\dfrac{\sqrt2m}{1+m^2}\)，\(y=-\dfrac{\sqrt2m^2}{1+m^2}\)，\(-1<m<1\).
\end{enumerate}
\end{solution}

\begin{problem}
设函数 \(f(x)=|2x+1|+|x-1|\).

\begin{enumerate}
\item 画出 \(y=f(x)\) 的图象；
\item 当 \(x\in[0,+\infty)\) 时，\(f(x)\le ax+b\)，求 \(a+b\) 的最小值.
\end{enumerate}

\bitmapfigure[max width=0.44\linewidth]{img/2018/national_paper_3_science/q23_fig1.png}
\end{problem}

\begin{solution}
\begin{enumerate}
\item 按 \(2x+1=0\) 和 \(x-1=0\) 分段. 于是
\[
f(x)=
\begin{cases}
-3x, & x<-\dfrac12,\\
x+2, & -\dfrac12\le x\le1,\\
3x, & x>1.
\end{cases}
\]
因此图象由三条直线段或射线组成：当 \(x<-\dfrac12\) 时为直线 \(y=-3x\) 的左侧射线；当 \(-\dfrac12\le x\le1\) 时为直线 \(y=x+2\) 的线段；当 \(x>1\) 时为直线 \(y=3x\) 的右侧射线. 三段在点 \(\left(-\dfrac12,\dfrac32\right)\) 和 \(\left(1,3\right)\) 处连接，这就确定了所求图象.
\item 当 \(x\ge0\) 时，
\[
f(x)=
\begin{cases}
x+2, & 0\le x\le1,\\
3x, & x\ge1.
\end{cases}
\]
若 \(f(x)\le ax+b\) 对一切 \(x\ge0\) 成立，由 \(x\ge1\) 时 \(f(x)=3x\) 可知 \(a\ge3\). 当 \(0\le x\le1\) 时，不等式给出
\[
b\ge 2-(a-1)x
\]
右端在该区间的最大值为 \(2\)，所以 \(b\ge2\). 于是
\[
a+b\ge3+2=5
\]
取 \(a=3\)，\(b=2\) 时，\(0\le x\le1\) 有 \(x+2\le3x+2\)，\(x\ge1\) 有 \(3x\le3x+2\)，条件成立. 因此 \(a+b\) 的最小值为 \(5\).
\end{enumerate}
\end{solution}
